\documentclass[main]{subfiles}

\title[Advanced Hyperparameter Optimization]{Advanced Hyperparameter Optimization}
\subtitle{Warmstarting From Prior Observations}
\author[Katharina Eggensperger]{Katharina Eggensperger}
\date{}

\titlebackground*{images/AutoML_Speeding_up_HPO.jpg}

%\institute{}
%\date{}

\begin{document}
	
	\maketitle
	

%----------------------------------------------------------------------
%----------------------------------------------------------------------
\begin{frame}[c]{Why Should We Start From Scratch?}

In practice, we have more data available from:
\begin{itemize}
    \item tune the same model repeatedly
    \item previous work on similar datasets
    \item re-run experiments with slight modifications
\end{itemize}

\vspace{1cm}
\textbf{$\leadsto$ How can we use this information to increase the efficiency of Bayesian optimization?} \\

\end{frame}
%-----------------------------------------------------------------------
%----------------------------------------------------------------------
\begin{frame}{A Motivating Example}

\begin{figure}
    \centering
    \includegraphics[width=0.5\linewidth]{05_Advanced_HPO/images/lr_bs_rule_of-thumb.png}
    \caption{Source: \liturl{Li et al. NeurIPS'24}{https://proceedings.neurips.cc/paper_files/paper/2024/file/ef74413c7bf1d915c3e45c72e19a5d32-Paper-Conference.pdf}}
\end{figure}
\end{frame}

%-----------------------------------------------------------------------
%----------------------------------------------------------------------
\begin{frame}[c]{The Initial Design}
In Bayesian optimization, we iteratively build a surrogate model: 

\begin{equation*}
    \surro_N (\conf | \hist), \qquad \hist = D_{init} \cup \{(\conf_i, \cost_i)\}_{i=n}^N
\end{equation*}

\vspace{0.5cm}
We start with an initial design $\dataset_{init}$:
\begin{itemize}
    \item typically random or space-filling designs (e.g., Latin Hypercube)
    \item efficiency of early iterations is highly sensitive to the initial design.
\end{itemize}

\vspace{0.5cm}
\textbf{$\leadsto$ How do we inject knowledge into this uninformed initialization?}

\end{frame}
%----------------------------------------------------------------------
%----------------------------------------------------------------------

\begin{frame}[c]{Prior Observations on the Same Task}

Assuming a fixed search space and given a target dataset $\dataset_t$, how can existing observations from manual exploration?

\vspace{0.1cm}
\textbf{Why?} Existing observations cover (human) intuition and serve as good starting points. \\
\textbf{Goal:} Inject knowledge into the early stages of Bayesian Optimization. \\
\textbf{Method:} (Partially) replace initial design. \\

\pause
\setbeamercolor{block title}{use=normal text,fg=normal text.fg}%
\setbeamercolor{block body}{use=normal text,fg=normal text.fg}%
\begin{columns}[T,onlytextwidth]
  \column{0.48\textwidth}
  \begin{block}{\raggedright Impact}
    \begin{itemize}
      \item reduces initial uncertainty volume
      \item biases surrogate towards existing (good) observations
      \item accelerates early exploitation
    \end{itemize}
  \end{block}

  \column{0.48\textwidth}
  \begin{block}{\raggedright Limitations}
    \begin{itemize}
      \item does not work across tasks
      \item performance depends on prior observations
    \end{itemize}
  \end{block}
\end{columns}

\end{frame}
%-----------------------------------------------------------------------
%----------------------------------------------------------------------
\begin{frame}{Alternatives}

Based on existing observations from a manual exploration, we can also:

\begin{itemize}
    \item choose a better surrogate model (e.g., kernel or model family)
    \item learn hyperparameters of surrogate (e.g., lengthscales)
    \item estimate noise of the objective function
    \item adjust the search space
\end{itemize}

\end{frame}
%-----------------------------------------------------------------------
%----------------------------------------------------------------------
\begin{frame}{Conclusion and Key Takeaway}

Prior observations can increase the efficiency of Bayesian optimization, e.g., by
\begin{itemize}
    \item (partially) replacing the initial design with prior observations,
    \item restricting the search space,
    \item using stronger priors/surrogates (e.g., kernel hyperparameters and noise estimation)
\end{itemize}

\vspace{0.5cm}
However, if prior observations are inconsistent (e.g., due to changes in code/data/seed) or biased, they can mislead optimization.

\end{frame}

\end{document}
